\chapter{Conclusions}
\section{Summary}
\label{sec:conclusion_summary}

The goal of this thesis was to optimize an existing HDC model for EMG-based classification of foot movements with respect to both classification accuracy and hardware efficiency. The original model achieves high classification accuracy but relies on high-dimensional hypervectors and a comparatively large CCIM, resulting in considerable memory and computational requirements. Therefore, the influence of the main HDC parameters and the representation of continuous input values was investigated first. Based on these results, the optimized model was subsequently implemented as a pipelined SystemC accelerator and synthesized for FPGA hardware.

The analysis of the HDC representation showed that, although the influence of the vector dimension $D$ strongly depends on the individual dataset, $D$ can be reduced by a factor of approximately $9.8$. While reducing the vector dimension alone also decreases the classification accuracy, it was shown that optimizing the feature-specific CCIM representation with the genetic algorithm can compensate for this accuracy loss if $D$ is kept above an empirical stability threshold of $ D= 751$ for the dataset-average accuracy. At the hardware-aligned dimension \(D=1024\) above this threshold, the optimized model achieves a dataset-average accuracy of \(95.96\%\), compared with \(94.78\%\) for the original \(D=10{,}000\) reference configuration. In contrast, reducing the number of quantization levels while preserving the reference accuracy was not achieved with the investigated quantization methods. They reduce the sensitivity of classification accuracy to the number of levels $L$, particularly for configurations with a small number of levels, but do not provide a consistent advantage over the original uniform quantization at $L=40$. Consequently, the final hardware configuration retains $L=40$ and uniform quantization, while reducing the vector dimension to $D=1024$. This reduces not only the CCIM size from $1.6~\mathrm{MB}$ to $163.84~\mathrm{kB}$ by a factor of 9.8, but also the number of elements processed in dimension-wise vector operations by the same factor.

The resulting HDC model was implemented as a pipelined FPGA accelerator including spatial encoding, temporal N-gram encoding, class-vector accumulation during training, and Hamming distance computation during inference. Two implementations with different degrees of word-level parallelism were evaluated. The fully parallelized implementation processes all 16 hypervector words in parallel and reaches a steady-state inference throughput of $8.33~\mathrm{MSamples/s}$ at a frequency of $100~\mathrm{MHz}$. However, its resource requirements prevent implementation on the smaller XCZU3EG FPGA that is suitable for embedded deployment. Reducing the parallelism of spatial encoding and class-vector accumulation to four words at a time also allows the CCIM banking to be reduced substantially. On the same XCVU57P FPGA, this decreases the LUT requirement by $48.6~\%$ and the BRAM requirement by $78.9~\%$ compared to the fully parallel version. The resulting partially parallelized accelerator can be implemented successfully on the XCZU3EG while meeting the target clock frequency of $100~\mathrm{MHz}$ and maintaining a throughput of $3.125~\mathrm{MSamples/s}$. RTL simulation of both architectures confirms that the synthesized implementations reproduce the expected HDC results.

Even the lower throughput of the partially parallel accelerator remains far above the requirements of the EMG application. With samples arriving at an interval of $9~\mathrm{ms}$, only approximately $111.1$ new samples are generated per second, compared with the $3.125$ million samples per second supported by the final accelerator. The accelerator therefore provides a throughput margin of more than four orders of magnitude and operates far above the minimum performance required for real-time processing.

Overall, the results show that the hardware requirements of HDC can be lowered while maintaining the dataset-average classification accuracy of the reference configuration. In particular, the combination of an optimized CCIM representation with a reduced vector dimension provides a considerably more compact HDC model, while the subsequent adaptation of the word-level parallelism and memory banking enables its implementation on an FPGA for embedded use. At the same time, the remaining processing performance exceeds the required EMG input rate by several orders of magnitude. The results therefore show that optimization of the HDC representation and adaptation of the hardware architecture complement each other. Reducing the vector dimension decreases the amount of data that has to be stored and processed, while adapting the hardware parallelism allows the resulting accelerator to be matched to the available FPGA resources.



\section{Limitations and Future Work}
\label{sec:limitations}

The results of this work have to be considered within the scope of the available datasets. The classification experiments are based on four subject-specific EMG datasets. While these datasets allow the influence of the investigated HDC parameters to be evaluated, generalization to a larger population cannot be assured. In particular, the optimized CCIM representations are evaluated within the respective dataset, and the transferability of an optimized representation between different subjects is not investigated. Future work could therefore evaluate the optimized HDC representation on a larger number of subjects and investigate inter-subject generalization, in particular whether class-vector separability is a useful GA objective for improving transferability. A subject-independent or jointly optimized CCIM could reduce the need for individual optimization while providing a more general representation for EMG classification.

The FPGA evaluation is restricted to the final model configuration with $D=1024$, $L=40$, and $N=3$. While the algorithmic analysis covers a broader parameter space, the hardware results therefore do not show how different model configurations affect FPGA resource usage and performance. In addition, only two degrees of word-level parallelism are evaluated. These implementations demonstrate the resource-performance trade-off between a fully and a partially parallelized architecture, but do not provide a complete exploration of the possible hardware design space. Intermediate degrees of parallelism could therefore be investigated to determine the smallest degree of parallelism that still satisfies the application requirements while further reducing resource usage. The detailed stage measurements additionally show that spatial encoding determines the processing interval of the partially parallelized implementation, making this stage and its memory organization a relevant target for further hardware optimization.

Finally, the reported hardware performance describes the HDC accelerator and its AXI4-Stream interface, but not a complete embedded assistive system. EMG acquisition, RMS preprocessing, quantization, software control, and actuator response remain outside the evaluated hardware path. The results only show that the accelerator itself satisfies the real-time requirements, but do not provide an end-to-end latency evaluation of the complete system. Integrating the accelerator into a complete embedded system would allow this latency to be assessed under realistic operating conditions. In addition, the FPGA evaluation focuses on resource utilization and timing, while power consumption and energy per classification are not measured. These quantities should be included in future evaluations to assess the suitability of the accelerator for battery-powered wearable devices.
